\section{Proofs for Section~\ref{sec:zf}}\label{app:zf}
% Appendix to Section 5: proofs for the ZF schemas and PA induction, and the computational evidence.

This appendix gives the proofs for \Cref{sec:zf}, in the order of that section, the theory of pattern templates with
several metavariables (\Cref{thm:zf:multi}), two results on induction that the main text summarizes
(\Cref{prop:zf:indgen,prop:zf:sub}), and the computational evidence (\Cref{app:zf:evidence}). Proofs are adapted from
the research records named in the source tags and were re-checked in transcription.

\subsection{Preliminaries}\label{app:zf:prelim}

For a set $D$ of sentences, $C(D)$ is the set of positions at which all data have the same symbol, with all ancestors
also in $C(D)$; a \emph{slot} of $D$ is a position outside $C(D)$ whose parent is in $C(D)$. For a template $T$,
$\mathrm{skel}(T)$ is the set of positions outside metavariable occurrences and their arguments.

\begin{lemma}[Rigid prefix]\label{lem:zf:rigid}
\status{proved}\src{prior induction C2(a); single Lemma P}
Let $T\in\SO$ and $D\subseteq\inst(T)$. Every datum has, at every position of $\mathrm{skel}(T)$, the symbol that $T$
has there. Hence $\mathrm{skel}(T)\subseteq C(D)$, and every occurrence of a metavariable of $T$ lies in $C(D)$ or is a
slot of $D$. (This is the $\SO$ form of \Cref{lem:single:basic}(c).)
\end{lemma}
\begin{proof}
In $\SO$ all occurrences are rigid and have metavariable-free arguments, so an instance replaces each occurrence by a
plugged body and leaves $\mathrm{skel}(T)$ unchanged. The parent of an occurrence lies in $\mathrm{skel}(T)$.
\end{proof}

Over $L_\in$ with parameters, a metavariable-free argument term is a variable or a parameter, i.e.\ a leaf. Hence
plugging arguments into the holes of a body $\beta$ does not change its tree shape, and every non-hole position of
$\beta$ keeps its symbol in $\beta[\bar u]$. We read bound variables in locally nameless style, so that ``$v$ is bound
above $c$'' and ``$v$ is bound inside the subtree at $c$'' are well defined. $\top:=\forall w(w=w)$ and $\bot:=\neg\top$
are closed.

% ---------------------------------------------------------------------------------------------------------------
\subsection{Encodings: proofs for Section~\ref{sec:zf:fo}}\label{app:zf:enc}

\begin{proof}[Proof of \Cref{thm:zf:classify}]
(a) Every occurrence in \Cref{tab:zf:forms} applies $P$ to distinct variables bound above it.

(b) ``If''. Let $N:=N_1=\dots=N_m$ and $\sigma:=F[Z,\dots,Z]$ with the guards $G=\{\mathrm{fresh}(v,Z):v$ frame-bound,
$v\notin N\}$. Let $\sigma\theta$ satisfy $G$, with $\psi:=\theta(Z)$. The frame-bound names free in $\psi$ lie in $N$;
its other free names are parameters. Let $\beta$ be $\psi$ with its free occurrences of the names in $N$ replaced by
the holes, in argument order. At every occurrence the arguments are the same names, each bound there by the binder
that the template intends, so $\beta[\bar a_k]=\psi$ for all $k$ and $\sigma\theta=T^*[P:=\lambda\bar z.\beta]$.
Conversely $T^*[P:=\lambda\bar z.\beta]=F[\psi,\dots,\psi]$ with $\psi=\beta[\bar a_1]$, whose free frame-bound names
lie in $N$, so it satisfies $G$.

``Only if''. Suppose $N_k\ne N_l$ and $\inst(\sigma,G)=\inst(T^*)$ for a first-order schema $\sigma$, all of whose
metavariables occur, and freshness guards $G$. Let $s_1,s_2$ be the instances with bodies
$\varphi_1:=\bigwedge_{i\le n}(z_i=z_i)$ and $\varphi_2:=\neg\varphi_1$, and $L:=\lgg(s_1,s_2)$. Frame positions are
common. At each $P$-position the roots differ, so $L$ has a metavariable there. The columns at occurrences $k$ and $l$
differ, because $\varphi_1$ mentions every argument and the tuples differ at some place. So $L$ has distinct
metavariables $Z_k\ne Z_l$ there. Since $s_1,s_2\in\inst(\sigma)$, $\sigma\gen L$, say $L=\sigma\rho$. Let $s_3$ be $L$
with $Z_k:=\top$ and every other metavariable $:=\bot$; then $s_3=\sigma(\rho\theta_3)$. For each metavariable $Y$ of
$\sigma$, the free names of $\rho\theta_3(Y)$ are those of the rigid part of $\rho(Y)$, which are also free in
$\rho\theta_1(Y)$, the value of $Y$ in $s_1$ (first-order matching is unique). Freshness guards are preserved when free
names are removed, and $s_1$ satisfies $G$; so $s_3\in\inst(\sigma,G)$. But $s_3\notin\inst(T^*)$. If $\beta[\bar a_k]$
is the closed $\top$, then $\beta$ has no hole, since a hole would leave the variable $a_{k,i}$ free in the slot; so
$\beta=\top$, and occurrence $l$ would also read $\top$, not $\bot$.

(c) The same proof with indices in place of names. Loose-index guards are also preserved when loose indices are removed.
At every occurrence the loose indices of the value must be smaller than its binder depth, so well-formedness implies
$\mathrm{loose}(Z)\subseteq\{0,\dots,d-1\}$ and nothing more. Since $I_1=I_k$ for all $k$, $I_1\subseteq\{0,\dots,d-1\}$;
so the guard $\mathrm{loose}(Z)\subseteq I_1$ is implied iff $I_1=\{0,\dots,d-1\}$.
\end{proof}

For \Cref{tab:zf:classify}: in de Bruijn notation $\Sep$ reads $P(\#0,\#2)$ at depth 3, so the guard ``$\#1$ (that is,
$y$) not loose'' is needed. $\Sep_{\mathrm J}$ reads $P(\#0)$ at depth 3. All three occurrences of $\EInd$ read
$P(\#0)$, at depths 2, 1, 1. $\Coll$ and $\ReplU$ read $P(\#1,\#0,\#2)$ in the antecedent and $P(\#1,\#0,\#3)$ in the
consequent. $\ReplS$ reads $P(\#1,\#0,\#2)$, $P(\#2,\#0,\#3)$, $P(\#1,\#0,\#3)$. $\ReplJ$ reads $P(\#2,\#1)$,
$P(\#2,\#0)$, $P(\#0,\#1)$. With $\varphi(x,y)$ only, $\Coll$ reads $P(\#1,\#0)$ twice at depths 3 and 4. There $\#2$ is
$A$ in the antecedent and $Y$ in the consequent, so the guard $\mathrm{loose}(Z)\subseteq\{\#0,\#1\}$ is needed: without
it, $Z:=(\#0=\#2)$ reads $y=A$ in the antecedent and $y=Y$ in the consequent, and the instance forces $Y\in Y$.

\begin{table}[t]
\centering\small
\begin{tabularx}{\textwidth}{@{}l >{\raggedright\arraybackslash}X >{\raggedright\arraybackslash}X c@{}}
\toprule
schema & named first-order (textbook names) & de Bruijn first-order & named, $\alpha$-varied\\
\midrule
$\Sep$ & yes, guard $\mathrm{fresh}(y,Z)$ & yes, guard ``$\#1$ not loose'' (not implied) & yes\\
$\Sep_{\mathrm J}$ & yes, guards $\mathrm{fresh}(X,Z)$, $\mathrm{fresh}(Y,Z)$ & yes, guard
  $\mathrm{loose}(Z)\subseteq\{\#0\}$ (not implied) & yes\\
$\EInd$ & \textbf{no} ($y$ vs.\ $x$) & \textbf{yes}, guard implied by well-formedness & no\\
$\Coll$, $\ReplU$ & yes, guard $\mathrm{fresh}(Y,Z)$ & \textbf{no} ($A$ is $\#2$ vs.\ $\#3$) & \textbf{no}\\
$\ReplS$, $\ReplJ$ & no ($y$ vs.\ $u$) & no & no\\
$\Coll$, $\ReplU$, $\varphi(x,y)$ only & yes, guards $\mathrm{fresh}(Y,Z)$, $\mathrm{fresh}(A,Z)$ & yes, guard
  $\mathrm{loose}(Z)\subseteq\{\#0,\#1\}$ (not implied) & no\\
\bottomrule
\end{tabularx}
\caption[Which ZF schemas are first-order patterns]{Is the schema a guarded first-order pattern? (By
\Cref{thm:zf:classify,prop:zf:alpha}; every schema is in $\PAT$.) Without its guard, every guarded entry has a false
instance (\Cref{prop:zf:false}).}
\label{tab:zf:classify}
\end{table}

\begin{proof}[Proof of \Cref{prop:zf:alpha}]
\src{cases Prop B1$\alpha$} The hypothesis class is first-order schemas over the named signature, with formula
metavariables and \emph{name metavariables} (at binder-name positions and variable leaves), and guards
$\mathrm{fresh}(v,Z)$ and $\mathrm{distinct}(v,w)$.
``If'': let $b_1,\dots,b_K$ be the frame's binders. Let $\sigma$ be the frame with the name of $b_k$ (at the binder and
at every leaf it binds) replaced by a name metavariable $v_k$, and every occurrence of $P$ replaced by one formula
metavariable $Z$. Its guards are $\mathrm{distinct}(v_k,v_l)$ for $k\ne l$, and $\mathrm{fresh}(v_k,Z)$ for every $b_k$
that is not an argument binder of $P$. An instance gives the binders distinct names and $Z$ a formula $\psi$ whose free
names are argument names or names the frame does not use. Let $\beta$ be $\psi$ with the argument names replaced by holes.
Every occurrence applies $P$ to the same binders, hence to the same names, so every slot reads $\beta[\text{names}]=\psi$,
and the instance is an $\alpha$-variant of $T^*[P:=\lambda\bar z.\beta]$. The converse is read off in the same way.
``Only if'': if occurrences $k,l$ differ at argument place $i$ (binders $b\ne b'$), fix one naming with all frame
binders distinct and argue as in \Cref{thm:zf:classify}. The slot texts at $k$ and $l$ differ in $s_1$ (one mentions the
name of $b$, the other that of $b'$), and distinct-guards only involve name positions, whose values in $s_3$ are those in
$s_1$. The ZF classification: $\Sep$ and $\Sep_{\mathrm J}$ have one occurrence. $\EInd$ applies $P$ to the binders
$\forall y$, $\forall x$ (first) and $\forall x$ (second). In $\Coll$ and $\ReplU$ the antecedent's $x,y$ and the
consequent's $x,y$ have different binders. $\ReplS$ and $\ReplJ$ differ even in names.
\end{proof}

\begin{proof}[Proof of \Cref{prop:zf:folgg}]
Frame positions are common columns. By \evR{} every $P$-position is a disagreement column, and the lgg merges equal
columns. So positions $k,l$ share a metavariable iff $\varphi_j[\bar a_k]=\varphi_j[\bar a_l]$ for all $j$. At a free
occurrence of $z_i$ the two texts carry $a_{k,i}$ and $a_{l,i}$, so the texts are equal iff the tuples agree at every
place used by $\varphi_j$.
For the last claim, let $U$ be the set of places used by some body. All positions share one metavariable $Z$, so for
$i\in U$ all occurrences carry the same name (index) $a_i$. The value of $Z$ in datum $j$ is $\psi_j=\varphi_j[\bar a_1]$,
whose free frame variables are exactly the $a_i$ with $z_i$ free in $\varphi_j$. So the most specific guards say that
the value's free frame variables lie in $\{a_i:i\in U\}$. For an instance $F[\psi,\dots,\psi]$ respecting them, let
$\beta$ be $\psi$ with each free $a_i$ ($i\in U$) replaced by $z_i$. Then $\beta[\bar a_k]=\psi$ for every $k$, so the
instance is $T^*[P:=\lambda\bar z.\beta]$.
\end{proof}

% ---------------------------------------------------------------------------------------------------------------
\subsection{First-order failures: proofs for Section~\ref{sec:zf:fofail}}\label{app:zf:fofail}

\paragraph{The false instances of \Cref{prop:zf:false}, complete list} \src{cases \S2.4}. Each is an instance of the
computed lgg of the two named data, respects its most specific guards, and is not an instance of the schema.
\begin{enumerate}[label=(\arabic*)]
\item $\EInd$, named; data $\EInd(x=x)$, $\EInd(\neg x\in a)$. The columns are $(y=y,\neg y\in a)$ at occurrence 1 and
  $(x=x,\neg x\in a)$ at occurrences 2 and 3, so the lgg is $\forall x(\forall y(y\in x\to Z_1)\to Z_2)\to\forall x Z_2$.
  Take $F_1$: $Z_1:=\neg y=y$, $Z_2:=\forall w\neg w\in x$. Its antecedent is $\forall x(\mathrm{empty}(x)\to
  \mathrm{empty}(x))$, which is valid, so $F_1$ is equivalent to the $\Pi_1$ sentence $\forall x\forall w\,\neg w\in x$.
  The $\Delta_0$ matrix of that sentence fails at the hereditarily finite sets $x=\{\emptyset\}$, $w=\emptyset$; by
  $\Delta_0$-absoluteness $F_1$ is false. ZF refutes it from ``some set has an element''.
\item $\ReplJ$, named; data $\ReplJ(x\in y)$, $\ReplJ(\neg x=a)$. The lgg is $\forall x\forall y\forall u(Z_1\wedge Z_2\to
  y=u)\to\forall X\exists Y\forall y(y\in Y\leftrightarrow\exists x(x\in X\wedge Z_1))$. Take $Z_1:=y=y$, $Z_2:=\neg u=u$.
  The antecedent is valid, and at $X=\{\emptyset\}$ the consequent gives a set $Y$ of all sets. Then $Y\in Y$, which
  contradicts Foundation; without Foundation, Separation yields Russell's set.
\item $\ReplJ$, de Bruijn: three blocks. $Z_1:=\bot$, $Z_2:=\bot$, $Z_3:=(y=y)$ give the same false consequent.
\item $\Coll$, de Bruijn; data $\Coll(y=x)$, $\Coll(\neg x=A)$. The two occurrences decouple. $Z_1:=y=y$, $Z_2:=\bot$
  give $\forall A(\forall x{\in}A\,\exists y\,y=y\to\exists Y\forall x{\in}A\,\exists y(y\in Y\wedge\bot))$, which is
  equivalent to $\forall A\neg\exists x\,x\in A$ and is refuted by the HF witness $A=\{\emptyset\}$.
\item $\ReplU$, de Bruijn; the same data. $Z_1:=y=x$ (and $\exists!y\,(y=x)$ is valid), $Z_2:=\bot$: equivalent to
  $\forall A\neg\exists x\,x\in A$.
\item $\ReplS$, de Bruijn; the same data, three blocks. $Z_1:=y=x$, $Z_2:=\bot$, $Z_3:=\bot$. The antecedent
  $\forall x{\in}A\,\exists y(y=x\wedge\forall u(\bot\to u=y))$ is valid, and the instance is equivalent to $\forall
  A\neg\exists x\,x\in A$.
\item Guard violations. $\Sep$ without $\mathrm{fresh}(y,Z)$: Russell's instance, in both encodings
  (\Cref{prop:zf:russell}). $\Coll$, $\ReplU$, $\ReplS$, named, without $\mathrm{fresh}(Y,Z)$: $Z:=(y=Y)$ gives
  \[\forall A\,(\forall x{\in}A\,\exists y\,(y=Y)\to\exists Y\,\forall x{\in}A\,\exists y\,(y\in Y\wedge y=Y)).\]
  Its antecedent
  holds, because $Y$ is a parameter there, and its consequent forces $Y\in Y$ for $A\ne\emptyset$, contradicting
  Foundation. (For $\ReplS$ add $\neg u=u$ at the uniqueness occurrence.) $\Coll$ with $\varphi(x,y)$ only, de Bruijn,
  without its guard: $Z:=(\#0=\#2)$, as in \Cref{app:zf:enc}.
\item Named, $\alpha$-varied. For $\Coll$, $\ReplU$ the blocks $(y=x)$, $\bot$, and for $\ReplS$ the blocks $(y=x)$,
  $\bot$, $\bot$, each give an instance equivalent to $\forall A\neg\exists x\,x\in A$.
\end{enumerate}

\paragraph{$\ReplS$, named, non-anchor data} \src{paper review, univzf\_repls\_pair.py; computed}. \Cref{prop:zf:coll} concerns data with
\evR{} and \zfN{y}. Without \evR, false instances arise when freshness guards are available only for formula
metavariables, the guard family used throughout \Cref{sec:zf}. The bodies $x\in y$ and $y\in A$ both
have root $\in$, so the named lgg of $\ReplS(x\in y)$ and $\ReplS(y\in A)$ is
\[\forall A\,(\forall x\,(x\in A\to\exists y\,(V_0\in V_1\wedge\forall u\,(V_2\in V_3\to u=y)))\to\exists Y\,\forall
x\,(x\in A\to\exists y\,(y\in Y\wedge V_0\in V_1))),\]
with term metavariables only, so no such guard is learned. Its instance $V_0:=Y$, $V_1:=y$, $V_2:=V_3:=u$ has $Y$ free
in the antecedent, i.e.\ a parameter there, and $y\in Y\wedge Y\in y$ in the consequent. At $A=\{\emptyset\}$ and
$Y=\emptyset$ the antecedent holds ($y:=\{\emptyset\}$; no set satisfies $u\in u$), while the consequent needs $y\in
Y'\in y$, which contradicts Foundation; so the instance is false and refutable in ZF. If freshness guards were also
learned for term metavariables ($\mathrm{fresh}(Y,V_0)$ holds on these data), this instance would be excluded; we have
not analysed that guard family.

\begin{lemma}[Deep leaf]\label{lem:zf:deepleaf}
\status{proved}\src{cases Lemma C1}
Let $s$ be a sentence of $L_\in$, $\ell$ a variable leaf of $s$, and $\tau$ a first-order schema with
$s\in\inst(\tau)$ and $|\tau|<\mathrm{depth}(\ell)$. Then $\tau$ has a formula metavariable $Z$ at a proper prefix $q$
of $\ell$ with $\mathrm{depth}(q)<|\tau|$. If $s|_q$ occurs in $s$ only at $q$, then $Z$ occurs in $\tau$ only at $q$, and
$s[q\leftarrow v]\in\inst(\tau)$ for every closed formula $v$. This remains true with freshness guards (named),
loose-index guards (de Bruijn), and distinct-guards on name metavariables ($\alpha$-varied named).
\end{lemma}
\begin{proof}
Follow the path to $\ell$ in $\tau$. Since $\tau$ has fewer than $\mathrm{depth}(\ell)$ nodes, the path leaves $\tau$ at
a leaf of $\tau$ at depth $<|\tau|$. That leaf is not a constant, because $s$ continues below it, so it is a
metavariable $Z$ at a proper prefix $q$ of $\ell$. In $L_\in$ every proper prefix of a variable leaf is a formula
position (binder-name children are not on the path), so $Z$ is a formula metavariable. If $Z$ also occurred at $q'$,
then $s|_{q'}=\theta(Z)=s|_q$, so $q'=q$. Changing $\theta(Z)$ to $v$ changes $s$ only at $q$. A closed $v$ has no free
names and no loose indices, and the other metavariables keep their values, so all guards still hold.
\end{proof}

\begin{proof}[Proof of \Cref{thm:zf:nounion}]
Let $\{\tau_1,\dots,\tau_r\}$ cover the schema. Take $n$ with $2n>\max_i|\tau_i|$, a genuine instance $s_n$ and a
designated leaf $\ell$ below; some $\tau_i$ covers $s_n$. By \Cref{lem:zf:deepleaf}, once we check that every subterm
along the path to $\ell$ occurs only once in $s_n$, $\tau_i$ contains $s_n[q\leftarrow v]$ for some proper prefix $q$
of $\ell$ and both $v\in\{\top,\bot\}$. It therefore suffices that, for every proper prefix $q$ of $\ell$, one of $\top,\bot$
makes $s_n[q\leftarrow v]$ false. ``Make the slot false'' means: if $q$ lies inside the slot below $e$ negations, choose
$v$ with $\neg^e v\equiv\bot$; the path inside the slot passes only through $\neg$, $\forall$ and $\wedge$, and
$\forall w\,v\equiv v$.

(i) $\EInd$, named. $\varphi_n(x):=\neg^{2n}\forall w\,\neg(w\in x)$ (``$x$ is empty''); $s_n=\EInd(\varphi_n)$; $\ell$ is
the $y$ of occurrence 1. For $q$ the root, take $v:=\bot$. For $q$ the antecedent or its $\forall x$-body, take
$v:=\top$: then $s_n$ is equivalent to $\forall x\,\varphi_n(x)$, ``every set is empty'', which is false. For $q$ the
formula $\forall y(y\in x\to\varphi_n(y))$ or its body, take $v:=\bot$: the antecedent becomes
$\forall x(\bot\to\varphi_n(x))$, true, and the conclusion is false. For $q$ inside $\varphi_n(y)$, make the slot false:
then $\forall y(y\in x\to\bot)$ says that $x$ is empty, the antecedent reads $\forall x(\mathrm{empty}(x)\to
\mathrm{empty}(x))$, which is true, and the conclusion is false. Every path subterm inside the slot contains $y$, which
occurs only in occurrence 1 and in the frame atom $y\in x$, so all path subterms are unique.

(ii) $\ReplJ$, both encodings. $\varphi_n(x,y):=\neg^{2n}(x\in y)$; $\ell$ is the $u$ of occurrence 2, $\varphi_n(x,u)$.
The consequent $C:=\forall X\exists Y\forall y(y\in Y\leftrightarrow\exists x{\in}X\,x\in y)$ is false: at
$X=\{\emptyset\}$, $Y$ would contain every $y$ with $\emptyset\in y$, and $y_0:=\{\emptyset,Y\}$ gives $Y\in y_0\in Y$,
contradicting Foundation. (Without Foundation, $R:=\{y\in Y:y\notin y\}$ and $y_1:=R\cup\{\emptyset\}$ give $y_1\in
y_1\leftrightarrow y_1\notin y_1$.) For $q$ the root take $\bot$. For $q$ the antecedent, its $\forall x,\forall
y,\forall u$ levels or its implication, take $\top$: the antecedent is true and $C$ is false. For $q$ the conjunction,
take $\bot$: $(\bot\to y=u)$ is true. Inside $\varphi_n(x,u)$, make the slot false. Uniqueness: the path subterms
contain $u$ (named), and in de Bruijn the designated occurrence's arguments $(\#2,\#0)$ differ from $(\#2,\#1)$ and
$(\#0,\#1)$.

(iii) $\Coll$, $\ReplU$, $\ReplS$, de Bruijn. $\varphi_n(x,y,A):=\neg^{2n}(y=x\wedge x\in A)$; $\ell$ is the $A$ in the
consequent's occurrence. The antecedent is true for every $A$, since $y:=x$ is the unique witness. Every proper prefix
$q$ is the root ($v:=\bot$), the $\forall A$-body ($v:=\bot$, giving $\forall A\,\bot$), or a subformula of the
consequent containing $\ell$. Above the slot, $v:=\bot$ turns the consequent into $\exists Y\bot$, $\exists Y\forall
x\bot$, or $\exists Y\forall x(x\in A\to\bot)$ (that is, ``$A$ is empty''); inside the slot (including the conjunction and
the atom $x\in A$), make the slot false. In each case the consequent is false for $A\ne\emptyset$, so
$\forall A(\text{true}\to\text{false for nonempty }A)$ is false. Uniqueness: in de Bruijn the slot's atom $x\in A$
reads $\#1\in\#3$, while the frame atoms read $\#0\in\#1$ (antecedent) and $\#0\in\#2$ (consequent), and the slot atom
$x\in A$ of the other occurrences reads $\#1\in\#2$ (antecedent occurrence) and, for $\ReplS$, $\#2\in\#3$ (uniqueness
occurrence). Every path subterm inside the slot contains the atom $\#1\in\#3$, and every path subterm above it contains
the slot; so all are unique. The argument tuple $(\#1,\#0,\#3)$ (resp.\ that of $\ReplS$'s third occurrence) differs
from the others.

(iv) $\Coll$, $\ReplU$, $\ReplS$, named under $\alpha$-variation. The same $\varphi_n$, with $s_n$ an $\alpha$-variant
of the instance in which all frame binders have distinct names. The designated leaf is the $y$ of the atom $y=x$ in the
consequent's occurrence. (The $A$-leaf of (iii) is not usable: in the named encoding the slot atom $x\in A$ also occurs as
a frame atom.) The case analysis is that of (iii); the only new prefix is the atom $y=x$ itself, where $v:=\bot$ makes
the slot false. Every path subterm inside the slot contains the consequent's own names for $x$ and $y$, which occur
nowhere outside the consequent, and inside the consequent the frame atoms are $x\in A$ and $y\in Y$, never $y=x$.
\end{proof}

\begin{proof}[Proof of \Cref{prop:zf:coll}]
(a) Let $L_S\theta$ satisfy the guards. The antecedent implies $\forall x{\in}A\,\exists y\,Z_1$ (drop the second
conjunct). $Z_1$'s value may mention $x,y,A$ and parameters, but not $Y$ (guard). The name $u$ is not bound at either
occurrence of $Z_1$, so a free $u$ would be a parameter at both. Collection, with these parameters, gives $\exists Y\forall
x{\in}A\,\exists y{\in}Y\,Z_1$, and ZF proves Collection. Strictness: $Z_1:=(y=x)$, $Z_2:=\bot$ respects the guards. It is
not an instance of $\ReplS$: a body $\beta$ with $\beta[x,y,A]=(y=x)$ is $z_2=z_1$, and then $\beta[x,u,A]=(u=x)\ne\bot$.
(b) ``If'' is the proof of (a), which uses only Collection. ``Only if'': let $\Coll(\varphi)$ be a Collection instance,
with $Y$ not free in $\varphi$. Renaming parameters gives an equivalent universal closure, so we may assume
$u\notin\FV(\varphi)$. Then $Z_1:=\varphi$, $Z_2:=\bot$ satisfy the guards, and the instance
$\forall A(\forall x(x\in A\to\exists y(\varphi\wedge\forall u(\bot\to u=y)))\to\exists Y\forall x(x\in A\to\exists
y(y\in Y\wedge\varphi)))$ is logically equivalent to $\Coll(\varphi)$, because $\forall u(\bot\to u=y)$ is valid.
\end{proof}

% ---------------------------------------------------------------------------------------------------------------
\subsection{Anchors for pattern targets: proofs for Section~\ref{sec:zf:pattern}}\label{app:zf:anchor}

\begin{proof}[Proof of \Cref{thm:zf:anchor}]
\emph{``If'', for $H=\SO$ (hence for all smaller $H$).} Let $T\in\SO$ cover $D$, with $D_j=T\eta_j$.

\emph{Step 1 (skeleton).} The root of $\varphi_j[\bar a_k]$ is $\mathrm{root}(\varphi_j)$. By \evR{} every slot root of
$T^*$ is a disagreement position, while frame positions are common. By \Cref{lem:zf:rigid}, $\mathrm{skel}(T)$ is a
prefix of the frame, and $T$'s occurrences lie at frame positions or at slot roots. For such a position $c$ put
$\kappa(c):=T^*|_c$, so that $D_j|_c=\kappa(c)[P:=\lambda\bar z.\varphi_j]$.

\emph{Step 2 (one body per metavariable).} Let $M$ have occurrences $M(\bar u^1)$ at $c_1$, \dots, $M(\bar u^r)$ at
$c_r$. The arguments are bound variables in scope or parameters, possibly repeated. Let $\beta_j:=\eta_j(M)$, so that
$\beta_j[\bar u^s]=D_j|_{c_s}$.

(2a) \emph{Shape.} Let $p$ be a slot root of $\kappa(c_1)$. Then $D_j|_{c_1}$ has the formula symbol
$\mathrm{root}(\varphi_j)$ at $p$, so $p$ is a non-hole position of $\beta_j$, and $D_j|_{c_s}$ also has
$\mathrm{root}(\varphi_j)$ at $p$. If $\kappa(c_s)$ had a frame symbol at $p$, all $\varphi_j$ would share it. If $p$
lay strictly inside a slot of $\kappa(c_s)$ with root $p'$, the non-leaf frame symbol of $\kappa(c_1)$ at $p'$ would be a
symbol of $\beta_j$, hence of $D_j|_{c_s}$ at $p'$, which is $\mathrm{root}(\varphi_j)$, for all $j$. Both contradict
\evR. By symmetry, $\kappa(c_1)$ and $\kappa(c_s)$ have slots at the same positions and the same symbols at non-leaf
frame positions. Frame leaves may differ.

(2b) \emph{Frame leaves.} Let $q$ be a frame leaf of $\kappa(c_1)$ carrying the variable $v$. If $v$ is bound inside
$\kappa(c_1)$, then $\beta_1|_q$ is that variable (holes become $\bar u^1$, which are bound above $c_1$ or are
parameters), so $\kappa(c_s)|_q$ is the variable bound by the corresponding binder of $\kappa(c_s)$; keep $v$. If $v$ is
bound above $c_1$, then $\beta_1|_q$ is a hole $h_m$ with $u^1_m=v$ ($\beta_1$ has no free bound variable, and a
parameter is not $v$); record $m$. By (2a), $q$ is a frame leaf of $\kappa(c_s)$ as well, and
$\kappa(c_s)|_q=D_1|_{c_s}|_q=u^s_m$.

(2c) \emph{Slot arguments.} Let $\kappa(c_1)$ have the slot $P(\bar a)$ at $p$ and $\kappa(c_s)$ the slot $P(\bar b)$
at $p$. For an argument place $i$, use \zfN{i} to pick $j_i$ and a free occurrence $q'$ of $z_i$ in $\varphi_{j_i}$. At
$p\cdot q'$, $D_{j_i}|_{c_1}$ carries $a_i$. If $a_i$ is bound inside $\kappa(c_1)$, then $\beta_{j_i}|_{p\cdot q'}$ is
that variable, so $b_i$ is the variable bound by the corresponding binder of $\kappa(c_s)$; keep $a_i$. Otherwise
$\beta_{j_i}|_{p\cdot q'}=h_m$ with $u^1_m=a_i$; record $m$. Since $D_{j_i}|_{c_s}|_p=\varphi_{j_i}[\bar b]$ carries
$b_i$ at $q'$, $b_i=u^s_m$.

Let $B_M$ be $\kappa(c_1)$ with every recorded leaf or argument replaced by its recorded hole. Every variable bound above
$c_1$ that occurs in $\kappa(c_1)$ does so at a frame leaf or as a slot argument, so $B_M$ is closed except for holes and
the metavariable $P$. $B_M[\bar u^1]=\kappa(c_1)$, and $B_M[\bar u^s]=\kappa(c_s)$ for all $s$ by (2a)--(2c). In
particular every recorded $u^s_m$ is a bound variable, even when $\bar u^s$ contains parameters or repetitions; any valid
choice of $\beta_1$ and $\beta_{j_i}$ works. (A term metavariable at a frame leaf is the case where $\kappa(c)$ is a
variable.)

\emph{Step 3.} Let $\sigma(M):=\lambda\bar h.B_M$ for every metavariable $M$ of $T$. Each occurrence $M(\bar u^s)$ becomes
$B_M[\bar u^s]=\kappa(c_s)$, and $T$'s rigid part is $F$'s, so $T\sigma=T^*$. Hence $T\gen T^*$ and
$\inst(T)\supseteq\inst(T^*)$.

\emph{``Only if'' (witnesses in $\PAT$).} If all roots are $f$, replace each $P(\bar a_k)$ by $\Phi_f(\bar a_k)$, where
$\Phi_\in(\bar a)=g_1(\bar a)\in g_2(\bar a)$ and $\Phi_=(\bar a)=g_1(\bar a)=g_2(\bar a)$ (term metavariables, whose
values are projections or parameters), $\Phi_\neg(\bar a)=\neg Q(\bar a)$, $\Phi_\circ(\bar a)=Q_1(\bar
a)\circ Q_2(\bar a)$ for binary $\circ$, and $\Phi_{\mathsf Q}(\bar a)=\mathsf Qw.R(\bar a,w)$ for
$\mathsf Q\in\{\forall,\exists,\exists!\}$. All occurrences are pattern occurrences. The template covers $D$, since atoms
of $L_\in$ have variables or parameters as arguments, and it misses $T^*[P:=\lambda\bar z.\psi]$ for any $\psi$ with
another root. If \zfN{i} fails, replace each $P(\bar a_k)$ by $P'(\bar a_k$ without its $i$-th entry$)$. This covers $D$
and misses $T^*[P:=\lambda\bar z.(z_i=z_i)]$.
\end{proof}

\paragraph{Several metavariables.} ZF needs one metavariable, but untagged mixtures behave like templates with several
(\Cref{sec:many}). Then a Plotkin-(D)-type event appears, and it depends on the class. Let $T^*$ be a pattern template
over $L_\in$ with a metavariable-free, parameter-free frame, metavariables $P,Q,\dots$ (arities $n_P,n_Q,\dots$),
every occurrence a pattern occurrence, and data $D_j=T^*[P:=\lambda\bar z.\varphi_{P,j},\dots]$. Let \zfR{P} and
\zfN{P,i} be \evR{} and \zfN{i} for the bodies of $P$. For a map $\pi:[n_P]\to[n_Q]\cup\mathrm{Par}$, let
$\varphi\circ\pi$ replace each hole $z_i$ by $z_{\pi(i)}$ (or by the parameter $\pi(i)$). The events are:
\zfDpat{PQ}: no bijection $\pi$ has $\varphi_{Q,j}=\varphi_{P,j}\circ\pi$ for all $j$;
\zfDdt{P\to Q}: no map $\pi$ has this property;
\zfDroot{PQ}: some $j$ has $\mathrm{root}(\varphi_{P,j})\ne\mathrm{root}(\varphi_{Q,j})$.
Then \zfDroot{PQ} implies \zfDdt{P\to Q}, which implies \zfDpat{PQ}. Under \zfN{P,i} for all $P,i$, \zfDdt{P\to Q}
for all ordered pairs is the event (D$^+$) of \Cref{cor:single:special}(b): a map $\bar u$ from the variables at an
occurrence $P(\bar a)$ to terms at an occurrence $Q(\bar b)$ with $\varphi_{Q,j}[\bar b]=(\varphi_{P,j}[\bar a])[\bar
u]$ for all $j$ sends every $a_i$, which occurs by \zfN{P,i}, to some $b_{i'}$ or a parameter, and so is a map $\pi$;
conversely every such $\pi$ gives such a $\bar u$. So the $\DT$ case of \Cref{thm:zf:multi} is also a case of \Cref{thm:single:anchor}.

\begin{theorem}[Several metavariables]\label{thm:zf:multi}
\status{proved}\src{cases Thm E$'$}
Assume \zfR{P} and all \zfN{P,i}, for every metavariable $P$. Then $D$ is an anchor in $\PAT$ iff \zfDpat{PQ} holds
for all $P\ne Q$. It is an anchor in $\DT$ iff \zfDdt{P\to Q} holds for all ordered pairs. In $\SO$, \zfDroot{} for
all pairs implies that $D$ is an anchor, which implies \zfDdt{} for all ordered pairs; the second implication cannot
be reversed. Without \zfR{P} or \zfN{P,i}, $D$ is an anchor in no class between $\PAT$ and $\SO$.
\end{theorem}
The proof of \Cref{thm:zf:anchor} works once each slot of $M$'s first occurrence faces a slot of the \emph{same}
metavariable at every other occurrence. If it faces a $Q$-slot instead, there is a fixed argument map from $P$'s
bodies to $Q$'s (\Cref{lem:zf:facing}), which the events exclude. Conversely, a violating map lets one rewrite every
$Q(\bar b)$ as an occurrence of $P$.

The three classes differ. Take $T_B:=\forall x\forall y(P(x,y)\to Q(x))$ with bodies $(x\in y,x\in x)$ and
$(\neg x=y,\neg x=x)$. Pattern anti-unification returns $T_B$, but $\forall x\forall y(P(x,y)\to P(x,x))\in\DT$ covers
the data and misses $\forall x\forall y(x\in y\to\neg x=x)$. Take $T_A:=\forall x\forall y(P(x,y)\to Q(x,y))$ with
bodies $(x\in y,y\in x)$ and $(\neg x=y,\neg x=x)$. Here $D$ is a $\DT$ anchor, but $\forall x\forall y(G(x,y,x)\to
G(x,x,y))\in\SO$ covers it and misses the same sentence. So, unlike the one-metavariable case, $\SO$ and $\DT$ anchors
differ even over $L_\in$: a metavariable without a pattern occurrence may choose different arguments on different
data. Whether \zfDroot{} is necessary in $\SO$ is open; the bounded searches of \Cref{app:zf:evidence} suggest:

\begin{conjecture}\label{conj:zf:root}
\status{conjecture}\src{cases \S2.9, \S6}
\zfDroot{} is not necessary for anchors in $\SO$: some data satisfying \zfR{P}, all \zfN{P,i} and all \zfDdt{P\to Q}
but violating \zfDroot{PQ} for some pair are anchors in $\SO$.
\end{conjecture}

\begin{lemma}[Facing]\label{lem:zf:facing}
\status{proved}\src{cases Thm E$'$}
In the setting of \Cref{thm:zf:multi}, let $T\in\SO$ cover $D$, let $M$ be a metavariable of $T$ with occurrences
$M(\bar u^s)$ at $c_s$ and bodies $\beta_j$, and assume \zfR{P} for all $P$.
(1) $\kappa(c_1)$ and $\kappa(c_s)$ have slots at the same positions and the same symbols at non-leaf frame positions.
(2) If $\kappa(c_1)$ has a $P$-slot $P(\bar a)$ at $p$ and $\kappa(c_s)$ a $Q$-slot $Q(\bar b)$ at $p$, then
$\mathrm{root}(\varphi_{P,j})=\mathrm{root}(\varphi_{Q,j})$ for all $j$. If moreover $\bar u^1$ is a pattern tuple and
all \zfN{P,i} hold, there is a map $\pi:[n_P]\to[n_Q]\cup\mathrm{Par}$, independent of $j$, with
$\varphi_{Q,j}=\varphi_{P,j}\circ\pi$ for all $j$. If also $\bar u^s$ is a pattern tuple and all \zfN{Q,i} hold, $\pi$ is a
bijection onto $[n_Q]$.
\end{lemma}
\begin{proof}
(1) is (2a) above, using \zfR{} for the metavariable at each slot; the root claim of (2) is its first sentence. Let
$\gamma_j:=\beta_j|_p$, so that $\varphi_{P,j}[\bar a]$ and $\varphi_{Q,j}[\bar b]$ are $\gamma_j$ with its holes
replaced by $\bar u^1$, resp.\ $\bar u^s$. Define $\pi(i)$ for $i\le n_P$: by \zfN{P,i} pick $j$ and a free occurrence
$q'$ of $z_i$ in $\varphi_{P,j}$; the leaf of $\varphi_{P,j}[\bar a]$ at $q'$ is $a_i$. ($\alpha$) If $a_i$ is bound
above $c_1$, then $\gamma_j|_{q'}=h_m$ with $u^1_m=a_i$, and $m$ is unique because $\bar u^1$ is a pattern tuple. The
leaf of $\varphi_{Q,j}[\bar b]$ at $q'$ is $u^s_m$, which is free relative to $p$; so it is some $b_{i'}$ or a
parameter, and we put $\pi(i):=i'$, resp.\ that parameter. ($\beta$) If $a_i$ is bound by a binder at relative position
$r$ between $c_1$ and $p$, then $\gamma_j|_{q'}$ is the variable bound at $r$, and the leaf of $\varphi_{Q,j}[\bar b]$ at
$q'$ is the variable bound at $r$ in $\kappa(c_s)$, which lies above $p$; so it is some $b_{i'}$, and $\pi(i):=i'$. In
both cases $\pi(i)$ depends only on $a_i$ (through $m$ or $r$), not on $j$ or $q'$. Now compare the two plugged texts leaf
by leaf through $\gamma_j$. Constants, parameters and variables bound inside $\gamma_j$ agree. A hole $h_m$ gives
$u^1_m$, a free leaf of the slot and hence some $a_i$, and $u^s_m=b_{\pi(i)}$ (or the parameter $\pi(i)$). A variable
bound between $c_1$ and $p$ gives $a_i$ and $b_{\pi(i)}$. Hence $\varphi_{Q,j}=\varphi_{P,j}\circ\pi$. If $\bar u^s$ is
also a pattern tuple, every $u^s_m$ is a bound variable, so $\pi$ maps into $[n_Q]$. It is injective: distinct $a_i$ give
distinct $m$ or distinct $r$, hence distinct variables, and the two cases give variables bound above, resp.\ below,
$c_s$. Every $b_{i'}$ free in some $\varphi_{Q,j}[\bar b]$ arises in this way, so under \zfN{Q,i} $\pi$ is onto.
\end{proof}

\begin{proof}[Proof of \Cref{thm:zf:multi}]
``If'' parts. Let $T$ in the class cover $D$. Step 1 of \Cref{thm:zf:anchor} holds verbatim, using \zfR{P} for every
$P$. For each metavariable $M$ of $T$ choose $c_1$ to be a pattern occurrence in $\PAT$ and $\DT$, and any occurrence in
$\SO$. Then every slot of $\kappa(c_1)$ faces, in each $\kappa(c_s)$, a slot of the \emph{same} metavariable. In $\SO$
this follows from \zfDroot{} and the root claim of \Cref{lem:zf:facing}; in $\DT$ from \zfDdt{} and the map; in $\PAT$
(where $\bar u^s$ is a pattern tuple too) from \zfDpat{} and the bijection. Steps 2b, 2c and 3 of
\Cref{thm:zf:anchor} then apply slot by slot, with \zfN{P,i} for the metavariable of each slot, and give $T\gen T^*$.

``Only if''. If \zfDpat{PQ} fails via a bijection $\pi$, let $T'$ be $T^*$ with every $Q(\bar b)$ replaced by
$P(b_{\pi(1)},\dots,b_{\pi(n_P)})$, a pattern occurrence. Then $T'[P:=\lambda\bar z.\varphi_{P,j}]=D_j$, because
$\varphi_{P,j}\circ\pi$ plugged with $\bar b$ is $\varphi_{Q,j}[\bar b]$. $T'$ misses every instance of $T^*$ in which
the $P$- and $Q$-bodies have different roots. So $D$ is not an anchor in $\PAT$ or in any larger class. If
\zfDdt{P\to Q} fails via $\pi$, the same construction with arguments $b_{\pi(i)}$ or the parameter $\pi(i)$ gives
$T'\in\DT$ ($P$ keeps its pattern occurrences, and the new occurrences have metavariable-free arguments), which covers $D$
and misses the same instances. In (c), an $\SO$ anchor is a $\DT$ anchor, hence satisfies \zfDdt{} by (b).
Non-reversibility is the example $T_A$ above. There \zfDpat{} and \zfDdt{} hold: the only map fitting
the first datum is the swap, which fails on the second. The template $\forall x\forall y(G(x,y,x)\to G(x,x,y))$ covers
the data with $G:=\lambda h_1h_2h_3.\,h_3\in h_2$, resp.\ $\lambda h_1h_2h_3.\neg h_1=h_2$. It misses $\forall x\forall
y(x\in y\to\neg x=x)$, because a body with $\beta[x,y,x]=(x\in y)$ is $h_1\in h_2$ or $h_3\in h_2$, so that
$\beta[x,x,y]\in\{x\in x,y\in x\}$. Without \zfR{P} or \zfN{P,i}, the witnesses of \Cref{thm:zf:anchor} applied to $P$
work in every class.
\end{proof}

\begin{proof}[Proof of \Cref{thm:zf:rates}]
By \Cref{thm:zf:anchor}, ``not an anchor'' is $\bigcup_fA_f\cup\bigcup_iB_i$, where $A_f$ says that all roots are $f$
and $B_i$ that $z_i$ is never free. In inclusion--exclusion over these events, every intersection containing two
distinct $A_f$ is empty. The intersection of $A_F$ (or of no $A$, $F={*}$) with $B_i$, $i\in I$, is the event that all
$N$ samples lie in $C_{F,I}$, of probability $\mu(C_{F,I})^N$, and its sign is $(-1)^{[F\ne{*}]+|I|+1}$. The upper
bound is the union bound.
\end{proof}

% ---------------------------------------------------------------------------------------------------------------
\subsection{PA induction: proofs for Section~\ref{sec:zf:pa}}\label{app:zf:pa}

Write $\Ind(\varphi)=\mathrm{imp}(U,V)$ with $U=\mathrm{and}(a,s)$, $s=\forall x\,W$, $W=\mathrm{imp}(b,c)$ and
$V=\forall x\,d$. The \emph{motive slots} $a,b,c,d$ hold $\varphi[0/x]$, $\varphi$, $\varphi[Sx/x]$, $\varphi$, and
the \emph{five-node frame} consists of the positions of $\mathrm{imp},\mathrm{and},s,W,V$. Two facts are used throughout
\src{prior induction C3, Lemma A}: (i) substituting terms for $x$ changes neither the root of a formula nor its
\emph{logical size} $\ell$ (the number of atoms, connectives and quantifiers); (ii) if $x$ is free in $\varphi$ and
$\varphi[u/x]=\varphi[v/x]$, then $u=v$. Also $St\ne x$, $St\ne0$ and $St\ne t$ for every term $t$. With $L=\ell(\varphi)\ge1$,
the slots have logical sizes $a,b,c,d{:}\ L$; $W{:}\ 2L+1$; $s{:}\ 2L+2$; $V{:}\ L+1$; $U{:}\ 3L+3$; and the whole
sentence $4L+5$. These six values are pairwise distinct.

\begin{proof}[Proof of \Cref{prop:zf:indtwo}]
\src{prior induction B1, B3}
(a) For $\varphi_1=(x+0=x)$ and $\varphi_2=(0+x=x)$ the $a$-column is $(0+0=0,0+0=0)$ and stays. The $b$-column
$(x+0=x,0+x=x)$ gives $z_{(x,0)}+z_{(0,x)}=x$. The $c$-column $(Sx+0=Sx,0+Sx=Sx)$ gives $z_{(Sx,0)}+z_{(0,Sx)}=Sx$. The
$d$-column equals the $b$-column. The four index tuples are distinct. The displayed instance ($z_1=z_2=z_3:=0$,
$z_4:=S0$) is false: $0+0=x$ forces $x=0$, and then $0+S0=1=Sx$, so the premise holds, while $\forall x(0+0=x)$ fails
at $x=1$. A single instance $I$ is covered by the ground schema $I$. For $\Ind(x=x),\Ind(Sx=Sx)$ the lgg is
$(t=t)\wedge\forall x(u=u\to Su=Su)\to\forall x(u=u)$, and every sentence instance has the valid conclusion
$\forall x(u=u)$.
(b) By fact (i) the $a,b,c$ columns all have mixed roots, so each becomes a metavariable indexed by its column. If some
$\varphi_j$ has $x$ free, then $\varphi_j[0/x]$, $\varphi_j$, $\varphi_j[Sx/x]$ are pairwise distinct, so the three
columns and their metavariables are distinct; the $d$-column equals the $b$-column. The displayed instance has a true
premise and the false conclusion $\forall x(0=S0)$. The lgg differs from $\zfLinf$ only if all roots are equal
(probability $\sum_fp_f^N\le p_{\max}^{N-1}$) or no motive has $x$ free (probability $(1-q)^N$).
\end{proof}

\begin{proof}[Proof of \Cref{thm:zf:indunion}]
\src{prior induction B2}
Let $\varphi_n:=(g_n(x)=x)$ and $n<m$, $d:=m-n$, so that $g_m(t)=g_n(g_d(t))$ as terms. In the $b$-slot both left sides
begin with $n$ copies of $0+\cdot$, so the lgg descends $n$ times and meets the column $(x,g_d(x))$. Its roots are a
variable and $+$, giving $b:=z_{(x,g_d(x))}$, so the slot is $g_n(b)=x$. Likewise the $a$-slot meets $(0,g_d(0))$ and
gives $g_n(a)=0$, and the $c$-slot meets $(Sx,g_d(Sx))$ and gives $g_n(c)=Sx$. The $d$-slot repeats $b$. The tuples
indexing $a,b,c$ have first entries $0,x,Sx$, so the metavariables are distinct, and the lgg is $L_n$. (For any subset
of $F'$ with least index $n$ the same descent gives $L_n$.)
For $M\ge n$, $I^*_M\in\inst(L_n)$: put $a,b:=g_{M-n}(0)$ and $c:=g_{M-n}(S0)$, and use $g_n(g_{M-n}(t))=g_M(t)$.
$I^*_M$ is false: $g_M(0)$ denotes $0$ and $g_M(S0)$ denotes $1$; the base holds, $g_M(0)=x$ forces $x=0$ and then
$g_M(S0)=Sx$, while $\forall x(g_M(0)=x)$ fails at $1$. It is refuted in $\Q$ from the closed facts $g_M(0)=0$ and
$g_M(S0)=S(g_M(0))$, which $\Q$ proves by $\Sigma_1$-completeness, and $\neg(g_M(0)=S0)$. Assume $g_M(0)=x$; by
congruence and transitivity $g_M(S0)=Sx$. Discharging and generalizing gives the step clause, modus ponens gives
$\forall x(g_M(0)=x)$, and instantiating at $S0$ contradicts $\neg(g_M(0)=S0)$.
Now a schema covering $\Ind(\varphi_n)$ and $\Ind(\varphi_m)$ is $\gen L_n$ and contains $I^*_n$. A union of $k$
truth-sound schemas covering $k+1$ members of $F'$ would put two of them into one schema. If a finite union $\bigcup\inst(\tau_l)$
contained $F'$, some $\tau_l$ would cover two members of $F'$ and hence contain the false, $\Q$-refutable sentence
$I^*_n$; so no such union lies between $F'$ and $\Th(\N)$, or between $F'$ and a consistent extension of $\Q$. The
motives of $F'$ are atomic, so the same holds for open induction.
\end{proof}
\emph{Guards} \src{prior induction B2 remark}. Data from $F'$ instantiate $b$ and $c$ by terms containing $x$, and $a$
by closed terms, so the strongest occurrence and freshness guards learnable from $F'$ are $x\in\FV(b)$, $x\in\FV(c)$ and
``$a$ closed''. The padded instance $a:=0$, $b:=0\cdot x$, $c:=S0+0\cdot x$ satisfies them and has the truth values of
$I^*_n$.

\begin{proof}[Proof of \Cref{thm:zf:K0}]
\src{prior induction referee R1}
Idea: for each instance that is not trivially true, a convergent rewrite system over the numerals and a new constant
$g$ gives an algebra containing $\N$ in which the step clause fails at $g$; amalgamating these algebras over $\N$ gives
a model of all instances and all true closed quantifier-free sentences. In detail:
In a sentence instance, $z_1:=t$ is closed (it occurs outside the binder) and $z_2:=u$ has $\FV(u)\subseteq\{x\}$.
Let $k$ be the value of $t$, and let $\hat u(g)$ be $u$ with $x$ replaced by a new constant $g$ and each maximal closed
subterm replaced by its numeral.

\emph{Case T: $\hat u(x)$ is literally $\num k+x$.} In every $L_A$-algebra $M$ containing $\N$ as a subalgebra, $t+a=u(a)$
for every $a$, so the conclusion and hence the instance hold.

\emph{Case W: otherwise.} Let $s:=\num k+g$ and $r:=\hat u(g)$, so $s\ne r$. Consider ground terms over the numerals and
$g$, with the rules $R_\N$ that evaluate $+$ and $\cdot$ on numerals, and one more rule $\rho$: $\rho:=r\to s$ if $s$ is a
proper subterm of $r$, and $\rho:=s\to r$ otherwise. \emph{Termination}: if $r$ is a numeral, use the lexicographic
measure (number of $g$'s, number of closed non-numeral positions). If $\rho=s\to r$ with $|r|\ge|s|$, use (closed
non-numeral positions, occurrences of $s$); no $\rho$-step creates a new occurrence of $s$, since an occurrence straddling
the rewritten position would force $r=g$ or $r$ closed. Otherwise use (closed non-numeral positions, size). In the
first case a $\rho$-step removes a $g$ and an $R_\N$-step lowers the second component. In the other two cases an
$R_\N$-step lowers the first component, and a $\rho$-step keeps it (both sides contain $g$ and have no closed
non-numeral subterm, and the context above still contains $g$) and lowers the second. \emph{Confluence}: there are no critical pairs. The left sides of $R_\N$ are closed, with only numerals as
proper subterms; the left side of $\rho$ contains $g$, is $R_\N$-normal, and is ground. So $R_\N\cup\{\rho\}$ is convergent
(critical pair lemma and Newman's lemma). Let $M_{t,u}$ be the ground terms modulo convertibility. Distinct numerals are
distinct normal forms, so $\N$ is a subalgebra of $M_{t,u}$, and $M_{t,u}\models s=r$, i.e.\ $t+g=u(g)$. The step clause
fails at $g$. If $\rho=s\to r$, the term $\num k+Sg$ is normal, with root $+$, while the normal form of $Su(g)$ is $Sr$
or a numeral. If $\rho=r\to s$, the normal form of $Su(g)$ is $Ss$ or $s$, and $\num k+Sg$ is normal and differs from both.
So the instance holds in $M_{t,u}$, its antecedent being false.

\emph{Amalgamation.} Take one $M_{t,u}$ for each Case-W instance, rename them apart over $\N$, and let $M$ be the free
completion of their union: an operation whose arguments all lie in one $M_i$ is computed there (this is well defined on
$\N$, a subalgebra of each), and any other application is a new formal element. Each $M_i$ is a subalgebra of $M$.
Closed terms evaluate in $\N$, so $M$ satisfies every true closed quantifier-free sentence. Each Case-W instance keeps its
quantifier-free witness $g$, and each Case-T instance holds in every algebra containing $\N$. A refutation by sound rules
would make some false closed quantifier-free sentence true in $M$, but such sentences have the same truth value in $M$ as
in $\N$. For $\Delta_0$ premises with primitive $<$, interpret $a<b$ as ``$a,b\in\N$ and $a<b$''; bounded quantifiers
with standard bounds then range over standard elements, so $\Delta_0$ truth is preserved.
\end{proof}

\begin{proof}[Proof of \Cref{thm:zf:indanchor}]
\src{prior induction C3}
We write the proof for parameter-free motives, as in the earlier analysis. (With parameters, the statement for $\DT$
is \Cref{cor:single:special}(c) \src{single Cor D.3}; we also checked that the steps below go through unchanged.)
\emph{``If''.} Fix $\varphi^\circ\in D$ with $x$ free \evN{} and let $T\in\DT$ cover $D$.
\emph{Step 1.} By \evR{} and fact (i) the four motive slots have different roots in different data, so by
\Cref{lem:zf:rigid} $\mathrm{skel}(T)$ is a prefix of the five-node frame. The occurrences of $T$ sit at some of the
positions of $\Ind$, $U$, $V$, $s$, $a$, $W$, $b$, $c$, $d$, at most one per position.
\emph{Step 2.} Every metavariable has a pattern occurrence at one of these formula positions, so all metavariables are
formula-valued, and all arguments are metavariable-free terms. Only $x$ is in scope at $W,b,c,d$, and nothing at the
others. So every metavariable is 0-ary, or unary with a pattern occurrence $M(x)$ at $W$, $b$, $c$ or $d$.
\emph{Step 3.} Let $\kappa(\cdot)$ denote slot contents in $\Ind(\varphi^\circ)$. A 0-ary metavariable takes a closed
value, so it cannot occur at $W,b,c,d$, whose contents have $x$ free. It cannot occur at two of $\Ind,U,V,s,a$, whose
contents have pairwise different logical sizes; so it occurs exactly once. A unary $M$ with a pattern occurrence at $\pi$
has the value $\lambda x.\kappa(\pi)$, and each occurrence $M(t)$ at $\sigma$ needs $\kappa(\pi)[t/x]=\kappa(\sigma)$; by
fact (i), $\sigma$ lies in $\pi$'s logical-size class, so $\pi=\sigma=W$ or $\pi,\sigma\in\{a,b,c,d\}$. If $M$ has a pattern
occurrence at $b$ or $d$, take $\pi$ there, so $\kappa(\pi)=\varphi^\circ$; fact (ii) forces $t=x$ at $b,d$, $t=Sx$ at $c$,
and $t=0$ at $a$. Otherwise $M$'s pattern occurrences are at $W$ (and then $M$ occurs only there) or at $c$; in the latter
case $\kappa(c)[t/x]=\varphi^\circ[St/x]$, and occurrences at $b,d$ or $a$ would need $St=x$ or $St=0$, which is impossible.
So $M$ occurs only at $c$, as $M(x)$.
\emph{Step 4.} Map a 0-ary metavariable at slot $\rho$ to $T_{\Ind}|_\rho$; $M$ with a pattern occurrence at $b$ or $d$ to
$\lambda y.P(y)$; $M$ with pattern occurrences only at $c$ to $\lambda y.P(Sy)$; and $M_W$ to $\lambda y.(P(y)\to P(Sy))$.
By Step 3 every occurrence of $T$ becomes $T_{\Ind}$'s subformula at its position, so $T\sigma=T_{\Ind}$ and $T\gen T_{\Ind}$.
\emph{``Only if''.} If all motives are closed, $Z_1\to\forall x\,Z_2$ (size 4) covers $D$ and misses $\Ind(x=x)$, whose
conclusion is not $\forall x\,Z_2$ with $Z_2$ closed. If all roots are $f$, use $(\Phi_f\wedge Z_1)\to Z_2$, where the
filler of the $a$-slot is $\Phi_f=z_1\mathbin{R}z_2$ for an atomic root $R\in\{=,<\}$ ($z_i$ 0-ary term metavariables),
$\neg Z$, $Z'\circ Z''$, or $\mathsf Qy\,K(y)$ with $K(y)$ a pattern occurrence. Each has size $\le7$, covers every
$\Ind(\varphi)$ with root $f$ (because $\varphi[0/x]$ has root $f$), and misses $\Ind(\neg x=0)$ (if $f$ is atomic) or
$\Ind(x=x)$. These witnesses lie in $\DT_s$ for every $s\ge7$.
\end{proof}

The following result is summarized after \Cref{thm:zf:indanchor}; sizes are as defined there.

\begin{proposition}[Generalizations, lgg, size threshold]\label{prop:zf:indgen}
\status{proved}\src{prior induction C4, C6.2}
(a) Up to renaming, exactly 26 templates of $\DT$ contain $\inst(T_{\Ind})$. Each takes a prefix of the frame, puts
fresh metavariables at the cut slots, and partitions the motive slots present, where $P(0)$ and $P(Sx)$ join a class
only together with a $P(x)$ slot. For every anchor these 26 are exactly the covering templates, so
$\lgg_{\DT}(D)=T_{\Ind}$.
(b) $\inst(T_{\Ind})$ is the intersection of the members of $\DT_s$ that contain it iff $s\ge12$. For $7\le s\le11$,
the cautious verifier over $\DT_s$ accepts on every anchor the false sentence
$(0=0\wedge\forall x(x=x\to Sx=Sx))\to\forall x(x=0)$.
\end{proposition}
\Cref{prop:zf:indpat} reads $T_{\mathrm{pat}}$ off the list in (a): it is the least generalization without
non-pattern occurrences. Part (b) is \Cref{lem:setting:closed} at work.

\begin{proof}[Proof of \Cref{prop:zf:indgen}]
\src{prior induction C4, C6.2; referee C}
(a) Every frame generalization is $\gen T_{\Ind}$ by the substitution of Step 4 above, so it contains $\inst(T_{\Ind})$.
Conversely, $\inst(T_{\Ind})$ contains anchors, and Steps 1--3 show that every template of $\DT$ covering an anchor is a
frame generalization. The count: the prefixes of the five-node frame are $\emptyset$, $\{\mathrm{imp}\}$, $\{\mathrm{imp},U\}$,
$\{\mathrm{imp},V\}$, $\{\mathrm{imp},U,V\}$, $\{\mathrm{imp},U,s\}$, $\{\mathrm{imp},U,s,V\}$, $\{\mathrm{imp},U,s,W\}$ and the
full frame. The motive slots present are none, none, $a$, $d$, $\{a,d\}$, $a$, $\{a,d\}$, $\{a,b,c\}$ and $\{a,b,c,d\}$.
The admissible partitions number $1,1,1,1,2,1,2,4$ and $13$ (for the full frame: $4$ with $b,d$ together, $9$ with $b,d$
apart), 26 in all. For an anchor, the covering templates are these 26, and $T_{\Ind}$ (full frame, one class) is below all
of them.
(b) Sizes count symbols, an occurrence $M(\bar t)$ counting $1+\sum|t_i|$. ``If'': $T_8=(P(0)\wedge Z)\to\forall
xP(x)$ (size 8), $T_{11}=(Z\wedge\forall x(P(x)\to P(Sx)))\to Z'$ (size 11) and $T_{12}=T_{\mathrm{pat}}$ (size 12)
contain $\inst(T_{\Ind})$. A sentence in $\inst(T_{12})$ has the form $\alpha\wedge\forall x(\beta\to\gamma)\to\forall
x\beta$ with $\beta$ a motive; $T_{11}$ adds $\gamma=\beta[Sx/x]$ and $T_8$ adds $\alpha=\beta[0/x]$. So the intersection
is $\inst(T_{\Ind})$. ``Only if'' ($s\le11$): sharing a metavariable between $d$ and $b$ needs $W$ and $V$ rigid (five
rigid nodes) and fillers of sizes $\ge1,2,\ge2,2$, so size $\ge12$; sharing between $d$ and $c$ needs $\ge13$. So templates
of size $\le11$ impose at most $a=d[0/x]$, $a=b[0/x]$, $c=b[Sx/x]$ and a frame prefix. The sentence
$s^*=(0=0\wedge\forall x(x=x\to Sx=Sx))\to\forall x(x=0)$ satisfies all three equations, but $b\ne d$; so it is no
induction instance, and it is false. Combine with \Cref{lem:setting:closed}.
\end{proof}

\begin{proof}[Proof of \Cref{prop:zf:indpat}]
\src{prior induction C9}
$\PAT\subseteq\DT$, and $D$ is an anchor in $\DT$ (\Cref{thm:zf:indanchor}). So the covering pattern templates are the
frame generalizations of \Cref{prop:zf:indgen}(a) without non-pattern occurrences: the $a$-slot gets a 0-ary
metavariable or is cut, and the $c$-slot gets its own $M_c(x)$ or is cut. The most specific of them is the full frame
with $b,d$ sharing $P(x)$, which is $T_{\mathrm{pat}}$; all others are above it. The displayed instance ($Z:=0=0$,
$P,Q:=\lambda x.\,x=0$) has a true premise and a false conclusion, and its $c$-slot $x=0$ is not $P(Sx)$.
\end{proof}

\begin{proof}[Proof of \Cref{prop:zf:indso}]
\src{prior induction C7}
``Only if''. If \evR{} fails, use the root-specialized templates above. If \zfB{} fails, let $\Theta=\{\theta_1,\dots,
\theta_k\}$ be the parent terms of the free occurrences of $x$ in the motives of $D$; each is $\ne x$ with
$\FV\subseteq\{x\}$. $T_\Theta:=Z\to\forall x\,P(\theta_1,\dots,\theta_k)$ covers $D$: in a motive, replace each maximal
occurrence of some $\theta_i$ in which $x$ is free by the hole $z_i$ (closed motives are their own bodies). It misses
$\Ind(x=x)$, because in its instances every free $x$ of the conclusion's body lies inside some plugged $\theta_i\ne x$.
``If''. Let $\varphi^\circ\in D$ have an unshielded free occurrence of $x$ at position $q$, and let $T\in\SO$ cover $D$. By
\evR{} and \Cref{lem:zf:rigid} the skeleton is a frame prefix, and all metavariables are formula-valued. For each $M$ fix a
body $\beta^\circ_M$ solving all its occurrences in $\Ind(\varphi^\circ)$. Plugging terms preserves logical size, so the
occurrences of $M$ lie in one class: $\{a,b,c,d\}$, $\{W\}$, or a single closed depth-0 slot.
\emph{Key fact.} If $M$ occurs at $b$ or $d$ with arguments $\bar t$, then $\beta^\circ_M$ has a hole $z_m$ exactly at $q$,
with $t_m=x$. The $x$ at $q$ must come from a plugged argument at some position $p\le q$; if $p<q$, the subterm at $p$ would
be a proper term $\ne x$ with $\FV\subseteq\{x\}$ containing the occurrence, contradicting unshieldedness. The same holds
at $W$, at position $1\cdot q$. At $c$, position $q$ holds $Sx$, so $\beta^\circ_M|_q$ is $z_m$ with $t_m=Sx$ or $S(z_m)$
with $t_m=x$; the hole cannot sit higher, because $q$'s parent is a predicate (a formula position) or a function symbol
whose subterm contains a variable bound inside $\varphi^\circ$, which no argument contains. At $a$, position $q$ holds $0$:
either a hole with argument $0$ or the constant $0$ of the body.
\emph{Substitution.} (i) $M$ occurring at $b$ or $d$, with hole $z_m$ at $q$: $M\mapsto\lambda\bar z.P(z_m)$. Every
occurrence $M(\bar u)$ at a motive slot has $u_m=\kappa(\cdot)|_q\in\{0,x,Sx\}$, which is $T_{\Ind}$'s argument there.
(ii) $M$ at $c$ but not at $b,d$: with $z_m$ at $q$, $M\mapsto\lambda\bar z.P(z_m)$ (arguments $0$ at $a$, $Sx$ at $c$);
with $S(z_m)$ at $q$, $M\mapsto\lambda\bar z.P(Sz_m)$, and then $M$ does not occur at $a$, since $S(u_m)=0$ is impossible.
(iii) $M$ only at $W$: $M\mapsto\lambda\bar z.(P(z_m)\to P(Sz_m))$. (iv) $M$ only at $a$: $M\mapsto\lambda\bar z.P(0)$.
(v) $M$ only at a depth-0 slot $\rho$: $M\mapsto\lambda\bar z.T_{\Ind}|_\rho$. Then $T\sigma=T_{\Ind}$.
\end{proof}

The Sub encoding summarized in \Cref{sec:zf:pa}:

\begin{proposition}[Sub encoding]\label{prop:zf:sub}
\status{proved}\src{prior induction C5, P1, P2, P6}
Record each induction as the step from $\mathrm{Sub}(\varphi,x,0,\alpha)$ and $\mathrm{Sub}(\varphi,x,Sx,\beta)$ to
$\alpha\wedge\forall x(\varphi\to\beta)\to\forall x\varphi$, where the decidable judgment $\mathrm{Sub}(\varphi,x,t,\psi)$
says that $\psi=\varphi[t/x]$ \citep{claude2026whatfollows}. With $x$ fixed, this is a first-order pattern in
$P,\alpha,\beta$. The lgg of genuine data recovers it iff \evR{} and \evN, so the failure probability is that of
\Cref{thm:zf:rates} for one argument place: $\sum_f p_f^N+(1-q)^N-\sum_f r_f^N$, where $r_f$ is the probability of
root $f$ with $x$ not free. If the induction variable varies, the variables must also not all be equal (or must be
$\alpha$-normalized). With the decidable guard ``all Sub premises are true'', the cautious verifier accepts only genuine
inductions.
\end{proposition}

\begin{proof}[Proof of \Cref{prop:zf:sub}]
\src{prior induction C5, P1, P2, P6}
The metavariables $P,\alpha,\beta$ receive $\varphi$, $\varphi[0/x]$ and $\varphi[Sx/x]$. By fact (i) the three \evR{}
events of Plotkin's lemma coincide with \evR. The three \evD{} events coincide with \evN, because
$\varphi\ne\varphi[0/x]$ iff $\varphi\ne\varphi[Sx/x]$ iff $\varphi[0/x]\ne\varphi[Sx/x]$ iff $x$ is free in $\varphi$: at a
free occurrence the three formulas carry $x$, $0$ and $Sx$. The rate is inclusion--exclusion over ``all roots are $f$''
(disjoint) and ``all motives closed''. If the induction variable is a metavariable $X$, its column consists of the
induction variables, which adds the event that they are not all equal. Genuine data satisfy the guard ``all Sub premises
true'', so by the most specific guard rule the verifier accepts $\inst(\lgg D)$ intersected with the Sub-true steps. A
Sub-true instance of $\sigma_{\mathrm{Sub}}$ has $\alpha=P[0/x]$ and $\beta=P[Sx/x]$ with $P$ a formula, i.e.\ it is a
genuine induction.
\end{proof}

\begin{proof}[Proof of \Cref{prop:zf:indeq}]
\src{prior induction B6}
(a) In first-order logic with equality, $\forall x(x=t\to\varphi)\leftrightarrow\varphi[t/x]$ when $x\notin\FV(t)$. Apply
this with $t=0$, $y$, $Sy$; since $\FV(\varphi)\subseteq\{x\}$, $y$ is not free in $\varphi$, and the step clause becomes
$\forall y(\varphi[y/x]\to\varphi[Sy/x])$, an $\alpha$-variant of $\forall x(\varphi\to\varphi[Sx/x])$. A sentence
instance $\IndEq(\psi)$ forces $\FV(\psi)\subseteq\{x\}$, since the first conjunct would leave any other free variable
free. Under universal closure, $P:=(x=0\vee\neg x=y)$ gives a false instance: at $y=1$ the premises hold (the step's
hypothesis holds only at $y=0$, and its conclusion always holds), while $\forall x(x=0\vee\neg x=1)$ fails at $x=1$.
(b) The four occurrences of $P$ carry the same $\psi_j$ in each datum and all other positions agree, so the lgg is
$\IndEq$ applied to the anti-unifier of the $\psi$-column, computed once (equal columns get equal results). It is
$\IndEq(P)$ iff the $\psi_j$ do not all have the same root; \evD{} is vacuous with one metavariable.
(c) By (a) the instances give every parameter-free induction instance. For $\varphi(x,\bar p)$ apply parameter-free
induction to $\psi(x):=\forall\bar p[\varphi(0,\bar p)\wedge\forall z(\varphi(z,\bar p)\to\varphi(Sz,\bar
p))\to\varphi(x,\bar p)]$. Here $\psi(0)$ is trivial, $\psi(x)\to\psi(Sx)$ uses the step clause, and $\forall x\,\psi(x)$
gives induction for $\varphi$ at every $\bar p$. Conversely, PA proves every instance. The trick raises quantifier
complexity, which is why parameter-free $\mathrm{I}\Sigma_n^-$ is weaker than $\mathrm{I}\Sigma_n$
\citep{kaye1988parameter}.
\end{proof}

\paragraph{Recorded motives in Lean and Metamath} \src{prior induction C12, P8, P9; referees A, C}.
In Lean~4 (version 4.32.0, checked locally in the earlier analysis and by its referee), $\mathtt{Nat.rec}$ takes the motive
$\{\mathit{motive}:\mathtt{Nat}\to\mathtt{Sort}\ u\}$ as an \emph{implicit} argument. The $\mathtt{induction}$ tactic
elaborates to $\mathtt{Nat.recAux}\ (\mathtt{fun}\ n\Rightarrow\dots)$, so the kernel term records the motive as a
$\lambda$. For eliminators marked \texttt{elab\_as\_elim} the elaborator infers the motive by abstracting the
discriminants, using \texttt{kabstract}. For closed
$\forall$-form steps the recorded motive is the unique matcher of $T_{\Ind}$, so recording it changes neither anchors
nor rates. In CIC proof terms, though, it is recoverable only up to conversion. Writing the motive as a $\lambda$-redex
$\mathrm{app}(\lambda x.P,t)$, with trusted $\beta$, gives a first-order pattern in one metavariable whose lgg is exact
iff the motives' roots are not all equal. But a \emph{named} outer binder $\forall n$ captures a parameter $n$ in $P$:
$P:=(x=0\vee\neg x=n)$ yields, after $\beta$, a sentence with a true antecedent and the false consequent $\forall
n(n=0\vee\neg n=n)$. So the binder must be nameless or guarded by $n\notin\FV(P)$.

Metamath's set.mm states induction (\texttt{finds}, \texttt{nnind}) with implicit-substitution hypotheses and
distinct-variable conditions \citep{megill2019metamath}. Transcribed to PA, the premises are H1
$x=0\to(\varphi\leftrightarrow\psi)$, H2 $x=y\to(\varphi\leftrightarrow\chi)$, H3 $x=Sy\to(\varphi\leftrightarrow\theta)$,
H4 $\psi$ and H5 $\chi\to\theta$, and the conclusion is $\varphi$. This is a first-order pattern in
$\varphi,\psi,\chi,\theta$. A set of guards from $\{x\notin\FV(\psi),x\notin\FV(\chi),x\notin\FV(\theta),y\notin
\FV(\varphi),x\ne y\}$ makes every instance a derived rule of PA iff it contains $y\notin\FV(\varphi)$ and one of
$x\notin\FV(\chi)$, $x\notin\FV(\theta)$ \status{proved}.
\emph{Sufficiency}: fix a model $M$ of $\PA$ in which the universal closures of the premises hold, a valuation $v$ in
$M$, and let $X=\{a\in M:M\models\varphi(v[x{:=}a])\}$. H1 and H4 give $0\in X$. For $a\in X$ put $w=v[x{:=}a,y{:=}a]$
and $w'=v[x{:=}a{+}1,y{:=}a]$. Then $\varphi(w)$ holds since $y\notin\FV(\varphi)$, and H2 gives $\chi(w)$. Then H5 gives
$\theta(w')$, either at $w'$ (if $x\notin\FV(\chi)$) or at $w$ (if $x\notin\FV(\theta)$). H3 at $w'$ gives
$\varphi(w')$, so $a+1\in X$. $X$ is definable in $M$ from the parameters of $v$, so induction in $M$ gives $X=M$.
Hence in every model of $\PA$ the universal closures of the premises imply that of the conclusion, i.e.\ the rule is a
derived rule of $\PA$.
\emph{Necessity}: $\varphi=\neg x=SSy$, $\psi=\neg0=SSy$, $\chi=\neg y=SSy$, $\theta=\neg Sy=SSy$ violates only
$y\notin\FV(\varphi)$, has PA-provable premises, and concludes the false $\neg x=SSy$. The instance $\varphi=(x=0)$,
$\psi=(0=0)$, $\chi=\theta=(x=0)$ violates only the two $x$-guards. So set.mm's conditions, read as ``not free'', are sound
but not minimal. Learned from canonical uses, the most specific guards include the spurious $y\notin\FV(\psi)$, which only
a non-canonical use removes.

% ---------------------------------------------------------------------------------------------------------------
\subsection{Computational evidence}\label{app:zf:evidence}

Everything above is proved. The computations corroborate it and found no counterexample. Scripts and outputs are in
\texttt{research/tracks/cases/} and \texttt{research/prior/induction/}. Each record also has an independent referee
implementation that imports none of the author's code. Bounded searches are evidence, not proof.
\begin{itemize}
\item \emph{Encodings and first-order failures} \src{cases z1, z4, z5, z8; referee r2}. The criterion of
  \Cref{thm:zf:classify} was computed from the frames. The witness $s_3$ was checked for all 7 non-first-order cases.
  The block structure of \Cref{prop:zf:folgg} matched on all pool pairs with \evR{} and on 2417 random pairs. Every
  exhibited false sentence is an instance of the computed lgg and not of the schema. For $n\le10$ the path subterms in
  \Cref{thm:zf:nounion}(i)--(iii) are unique, and 2136 random small generalizations $\tau$ of $s_n$ contain both
  $s_n[q\leftarrow\top]$ and $s_n[q\leftarrow\bot]$. The pure-logic steps of the case analysis hold on all
  $\in$-structures of size $\le3$. The guard criterion of \Cref{prop:zf:russell}(b) matched on all pool pairs and on 1537
  random pairs.
\item \emph{Anchors, pattern lgg and rates} \src{cases z2, z3, z6, z7, z9; referee r3, r4}. For each of the seven schemas,
  all 120 pairs of a 16-body pool were checked with a reduced $\DT$ enumeration, which is complete for anchor-hood:
  anchor $\iff$ \evR{}$\wedge$\zfN{} on 120/120. Full $\DT$ and $\SO$ enumerations on the smaller frames, and bounded
  $\DT$/$\SO$ searches on all frames ($7\times20$ pairs), agree, with identical $\SO$ and $\DT$ verdicts. The referee's
  $\SO$ search, with repeated and parameter arguments, found 0 disagreements in 460 runs. An $n$-ary pattern
  anti-unifier with a permutation merge is exact iff \evR{}$\wedge$\zfN{} on all pairs and 400 random triples per
  schema; the referee's implementation agrees on 3500 data sets, and 10\,500 random instances of its lggs are genuine.
  We did not run the published implementation of \citet{baumgartner2017higher}. For several metavariables, the
  pattern-lgg criterion held on $3\times1500$ random data sets and the $\DT$ criterion on $2\times150$ pairs. Every
  $\SO$ pair with \zfDroot{} was an anchor, and none violating \zfDdt{} was. The few pairs with \zfDdt{} but not
  \zfDroot{} were $\SO$ anchors under the bounded search, which suggests \Cref{conj:zf:root}. The
  exact rates agree with Monte Carlo.
\item \emph{Induction} \src{prior induction b1--b4, c0--c5; referees B, C}. An exactly false lgg instance was found for
  2357 of the 2556 pairs of quantifier-free motives of size $\le5$; the smallest unsound pair has motive sizes $3+3$.
  $\lgg=L_n$ for all $0\le n<m\le10$. For $K_0$, a congruence-closure test on 2224 instances finds the step clause
  refutable over $\N$ exactly in the 12 Case-T instances. For all 351 pairs of a 27-motive pool, all $\SO$ templates of
  size $\le14$ (arity $\le2$, arguments of size $\le3$; 4.6 million covering templates) were classified: $\DT$ anchor
  $\iff$ \evR{}$\wedge$\evN{}, and $\SO$ anchor $\iff$ \evR{}$\wedge$\zfB, on 351/351. The referee found no
  disagreement on 176 further data sets. With bounded arity, though, some non-\zfB{} sets are $\SO$ anchors, since
  \Cref{prop:zf:indso} needs unbounded arity. The 26 generalizations and the threshold 12 were reproduced. The Sub-lgg
  and $\DT$ verdicts coincide on all pairs and 2925 triples. $\IndEq$ and $\Ind$ agree on 30\,000 random pairs of a
  motive and a finite structure.
\item \emph{Experiment E2} \src{experiments E2}: see \Cref{tab:exp:e2} and \Cref{sec:exp}.
\end{itemize}
